\begin{enumerate}
\item
\begin{enumerate}
\item $P(\N)=1$ et par $\sigma$-additivité, $P(\N)=\dsum_{n=0}^{\pinf} P(\ens{n})$. La série $\dsum_{n\in\N} P(\ens{n})$ converge donc son terme général tend vers 0.
\item Le reste de cette même série tend vers 0 donc $\dlim _{n \rightarrow+\infty} P(\{k\in\N,\ k \geqslant n\})$.
\end{enumerate}
\item $(a_n)$ est strictement décroissante de limite nulle, donc tous ses termes sont strictement positifs. \\
Analyse : nous devons avoir $P(\N)=1$ donc $\l=\dfrac1{a_0}$.\\
De plus,
\begin{align*}
P(\{n\}) & =P\big[\ens{k\in\N,\ k \geqslant n} \backslash\ens{k\in\N,\ k \geqslant n+1}\big] \\
& =P(\ens{k\in\N,\ k \geqslant n})-P(\ens{k\in\N,\ k \geqslant n+1})=\frac{a_n-a_{n+1}}{a_0}.
\end{align*}
Synthèse : Posons $\l=\dfrac1{a_0}$. Pour tout $n\in\N$, posons $p_n=\l(a_{n}-a_{n+1})$. Alors $p_n>0$ et $\dsum_{n=0}^{\pinf}p_n=\l(a_0-\dlim_{n\to\pinf} a_n)=1$ par télescopage. Il existe donc une probabilté $P$ telle que pour tout $n\in\N$, $P(\ens{n})=p_n$. De plus $P(\N)=\l a_0=1$. Enfin, si $n\in\N$,
\begin{align*}
P(\ens{k\in\N,\ k \geqslant n})&=\dsum_{k=n}^{\pinf}P(\ens{n})= \dsum_{k=n}^{\pinf}P(\ens{n})\l(a_{k}-a_{k+1})\\
&= \lambda a_n.
\end{align*}
Il existe donc une probabilité vérifiant la propriété voulue \ssi\ $\l=\dfrac1{a_0}$.
\end{enumerate}
